\subsubsection{Classical Solution and Query Complexity}\label{ex:classical-solution-and-query-complexity}

Now that we know the problem about the Deutsch's algorithm, let us ask a natural question: \emph{\textbf{how many times must a \underline{classical computer} query the oracle to determine whether the function is constant or balanced?}} To answer this question, we first need to introduce the notion of \textbf{query complexity}.

\highspace
Recall that the function is provided as a \textbf{black box} (\textbf{oracle}). We are \textbf{not allowed to inspect its implementation}. The only permitted operation is to evaluate the function at a chosen input. For example, we may ask the oracle to evaluate $f(0)$ or $f(1)$. Each evaluation of the oracle is called a \definition{Query}.

\highspace
Since interacting with the oracle is assumed to be the expensive part of the computation, we measure the efficiency of an algorithm by \textbf{counting the number of oracle queries it performs}, rather than the total number of elementary operations. This measure is called the \definition{Query Complexity} of the algorithm.

\highspace
\begin{flushleft}
    \textcolor{Green3}{\faIcon{question-circle} \textbf{Can one classical query be enough?}}
\end{flushleft}
Suppose we decide to evaluate the function only once. For example, we query the oracle with $f(0)$. There are two possible outcomes:
\begin{itemize}
    \item Case $f(0) = 0$. If the oracle returns $0$, there are still two functions consistent with this observation:
    \begin{center}
        \begin{tabular}{@{} l c c l @{}}
            \toprule
            \textbf{Function} & $f(0)$ & $f(1)$ & \textbf{Type} \\
            \midrule
            $f(x) = 0$ & 0 & 0 & Constant \\
            $f(x) = x$ & 0 & 1 & Balanced \\
            \bottomrule
        \end{tabular}
    \end{center}
    Both functions produce exactly the same answer for the first query. Therefore, after observing only $f(0) = 0$, we still cannot determine whether the function is constant or balanced.
    
    
    \item Case $f(0) = 1$. Similarly, if the oracle returns $1$, there are still two functions consistent with this observation:
    \begin{center}
        \begin{tabular}{@{} l c c l @{}}
            \toprule
            \textbf{Function} & $f(0)$ & $f(1)$ & \textbf{Type} \\
            \midrule
            $f(x) = 1$ & 1 & 1 & Constant \\
            $f(x) = \overline{x}$ & 1 & 0 & Balanced \\
            \bottomrule
        \end{tabular}
    \end{center}
    Again, one query leaves two possible candidates.
\end{itemize}
\textcolor{Red2}{\faIcon{exclamation-triangle} \textbf{Why one query is not enough?}} After a single query, we know only one of the two values: $f(0)$ or $f(1)$. The value of the other input remains unknown.

\highspace
Since both a constant function and a balanced function can produce exactly the same answer for the queried input, \textbf{one classical query can never distinguish the two cases with certainty}. This is true regardless of whether we choose to evaluate $f(0)$ or $f(1)$: one unknown output always remains.

\highspace
So, to remove the ambiguity, we must evaluate \textbf{both} possible inputs. In other words, we need to query the oracle \textbf{twice} to determine with certainty whether the function is constant or balanced. The classical strategy can be visualized as the following decision tree:

\begin{center}
    \begin{tikzpicture}[
        node distance=0.8cm and 0.5cm,
        every node/.style={
            draw,
            rounded corners=2pt,
            minimum width=1.8cm,
            minimum height=0.8cm,
            align=center
        },
        decision/.style={
            -{Latex[length=2mm]},
            line width=0.45pt
        }
    ]

        \node (q0) {Query $f(0)$};

        \node (f00) [below left=of q0] {$f(0)=0$};
        \node (f01) [below right=of q0] {$f(0)=1$};

        \node (q10) [below=of f00] {Query $f(1)$};
        \node (q11) [below=of f01] {Query $f(1)$};

        \node (f100) [below=1.0cm of q10, xshift=-1.15cm] {$f(1)=0$};
        \node (f101) [below=1.0cm of q10, xshift= 1.15cm] {$f(1)=1$};

        \node (f110) [below=1.0cm of q11, xshift=-1.15cm] {$f(1)=0$};
        \node (f111) [below=1.0cm of q11, xshift= 1.15cm] {$f(1)=1$};

        \node (c0) [below=of f100] {Constant};
        \node (b0) [below=of f101] {Balanced};
        \node (b1) [below=of f110] {Balanced};
        \node (c1) [below=of f111] {Constant};

        \draw[decision] (q0) -- (f00);
        \draw[decision] (q0) -- (f01);

        \draw[decision] (f00) -- (q10);
        \draw[decision] (f01) -- (q11);

        \draw[decision] (q10) -- (f100);
        \draw[decision] (q10) -- (f101);

        \draw[decision] (q11) -- (f110);
        \draw[decision] (q11) -- (f111);

        \draw[decision] (f100) -- (c0);
        \draw[decision] (f101) -- (b0);
        \draw[decision] (f110) -- (b1);
        \draw[decision] (f111) -- (c1);

    \end{tikzpicture}
\end{center}

\highspace
\textcolor{Red2}{\faIcon{exclamation-triangle} \textbf{Classical query complexity.}} We can now summarize the classical solution. A \textbf{deterministic classical algorithm} cannot distinguish between a constant and a balanced function after only one oracle query. It must evaluate both possible inputs $f(0)$ and $f(1)$ requiring \textbf{exactly two orcale queries}.

\highspace
Therefore, the \hl{deterministic classical query complexity of Deutsch's problem is 2.} The remarkable result of Deutsch's algorithm is that a \textbf{quantum computer} can \textbf{solve the same problem} with \textbf{only one query} to the quantum oracle. And now the natural question arises: \emph{\textbf{how can a quantum computer possibly do it with only one query?}} The answer is provided in the next section.